回到一般情形. 对任意 $(A, B)$-双模 $P$, 模论的常识 \cite[定理 6.6.5]{Li1} 给出伴随对
\begin{equation}\label{eqn:AB-adjunction}\begin{tikzcd}
	(\cdot) \dotimes{A} P : \cated{Mod}A \arrow[shift left, r] & \cated{Mod}B: \Hom_{\cated{Mod}B}(P, \cdot). \arrow[shift left, l]
\end{tikzcd}\end{equation}
定理 \ref{prop:Morita} 说明精确到同构, 左右两端的函子分别穷尽了 $\mathrm{Fct}^c(\cated{Mod}A, \cated{Mod}B)$ 和 $\mathrm{Fct}_c(\cated{Mod}B, \cated{Mod}A)$ 的所有对象.

基于例 \ref{eg:adjunction-monad}, 伴随对 \eqref{eqn:AB-adjunction} 在 $\cated{Mod}A$ 上确定单子 $T$, 在 $\cated{Mod}B$ 上确定余单子 $L$, 其一般描述如下. 考虑右 $A$-模 $N$ 和右 $B$-模 $M$:
\begin{equation}\label{eqn:bimod-adjunction}
	\begin{array}{|c|c|} \hline
		\text{单位}\; \eta_N & \text{余单位}\; \epsilon_M \\ \hline
		N \to \Hom_{\cated{Mod}B}\left( P, N \dotimes{A} P \right) & \Hom_{\cated{Mod}B}(P, M) \dotimes{A} P \to M \\
		x \mapsto \left[ p \mapsto x \otimes p \right] & \varphi \otimes p \mapsto \varphi(p) \\ \hline
	\end{array}
\end{equation}
